% ==============================================================================
% Secção 3: Excelência técnico-científica
% ==============================================================================

\subsection{Visão a longo prazo}
\begin{campotexto}
A visão deste projeto passa por integrar o desenvolvimento de diferentes subsistemas do veículo elétrico, em vez de os otimizar de forma isolada, visando ganhos simultâneos de eficiência energética, segurança, conforto e sustentabilidade. Pretende-se, assim, estabelecer as bases técnico-científicas para conceber os carros do futuro, alinhando a inovação tecnológica com a evolução social e regulamentar (enquadramento PESTEL).

No domínio do armazenamento de energia, prevê-se a superação progressiva das limitações das baterias de lítio. Com a redução dos custos iniciais de fabrico, as baterias de estado sólido permitirão eliminar eletrólitos líquidos inflamáveis, oferecendo maior densidade energética e menor degradação térmica. Em alternativa, as soluções de sódio poderão afirmar-se pelo custo reduzido e pela abundância do material, desde que seja compensada a sua menor densidade energética, devendo ambas as opções ser acompanhadas por estratégias responsáveis de reciclagem e descarte em fim de vida.

Ao nível da dinâmica veicular, projeta-se a substituição gradual das suspensões passivas de calibração fixa por sistemas ativos, recorrendo a atuadores elétricos, eletromagnéticos ou hidráulicos de geometria variável para controlo do ângulo de camber, em articulação com sensores e inteligência artificial. Esta conjugação viabilizará uma atuação preditiva e a otimização contínua do comportamento em estrada, à medida que sejam resolvidos o elevado consumo dos atuadores e a dificuldade de calibração para múltiplos cenários de condução.

Complementarmente, procura-se alargar a recuperação energética para além da travagem e desaceleração. A adoção de amortecedores regenerativos (eletromagnéticos, hidráulicos ou eletro-hidráulicos) permitirá converter em eletricidade as vibrações mecânicas atualmente desperdiçadas, exigindo investigação dedicada para ultrapassar a complexidade de acoplamento ao sistema elétrico do veículo.

Em suma, a excelência técnico-científica desta proposta reside no reconhecimento de que o progresso automóvel não depende apenas de uma nova fórmula química, mas da articulação coerente entre armazenamento avançado, controlo inteligente e aproveitamento das vibrações. Em etapas subsequentes, será fundamental aprofundar a viabilidade industrial destas soluções e quantificar o seu impacto real na autonomia e no desempenho global.
\end{campotexto}

\subsection{Avanços científicos e tecnológicos}
\begin{campotexto}
O desenvolvimento dos carros do futuro exige compreender as limitações das soluções automóveis atuais e identificar os avanços científicos e tecnológicos capazes de as superar. A crescente eletrificação do setor impõe requisitos exigentes de eficiência energética, segurança, conforto e sustentabilidade, sendo o desempenho global dos veículos elétricos fortemente condicionado pelas características e pelo peso das baterias, bem como pela resposta dos sistemas de suspensão e pela capacidade de aproveitamento energético. Neste contexto, o projeto analisa o estado da arte e as tecnologias emergentes em três domínios complementares:
\begin{itemize}
    \item \textbf{Armazenamento de energia:} atualmente, as baterias de lítio (LFP/NMC) constituem a solução dominante nos veículos elétricos, mas apresentam constrangimentos ao nível da estabilidade térmica, dos eletrólitos líquidos inflamáveis, da dependência de matérias-primas críticas e do descarte em fim de vida. Como alternativas emergentes destacam-se as baterias de estado sólido, que substituem o eletrólito líquido por um sólido para eliminar o risco de inflamabilidade, reduzir a degradação térmica e aumentar a densidade energética, enfrentando ainda elevados custos iniciais de fabrico; e as baterias de sódio, que tiram partido da abundância deste elemento para oferecer menor custo de produção e perspetivas favoráveis de sustentabilidade e reciclagem, embora apresentem menor densidade energética. Paralelamente à química das células, torna-se essencial integrar estratégias eficazes de reciclagem e gestão de fim de vida.
    \item \textbf{Suspensão e amortecimento:} a solução mais comum continua a ser a suspensão passiva, composta por molas e amortecedores hidráulicos com calibração fixa, o que impõe um compromisso permanente entre conforto e estabilidade, agravado pelo peso das baterias. Os avanços recentes apontam para arquiteturas adaptativas em tempo real, nomeadamente a suspensão ativa --- recorrendo a atuadores elétricos ou eletromagnéticos, bem como a soluções hidráulicas de geometria variável para controlo do ângulo de \textit{camber} da roda --- capaz de aumentar o conforto e a estabilidade, embora com elevado consumo energético. Em articulação, o uso de sensores e Inteligência Artificial (IA) permite processar dados em tempo real para otimizar continuamente a resposta do sistema e viabilizar uma atuação preditiva, persistindo o desafio da calibração para múltiplos cenários de condução.
    \item \textbf{Recuperação de energia:} no estado da arte presente, a regeneração energética limita-se à travagem e desaceleração, desaproveitando a energia mecânica produzida pelas oscilações verticais durante a marcha. Os amortecedores regenerativos surgem como tecnologia emergente ao utilizarem dispositivos eletromagnéticos, hidráulicos ou eletro-hidráulicos para converter parte dessas vibrações em eletricidade --- como em arquiteturas com unidade de engrenagens e alternador acopladas a baterias auxiliares de 48~V. Apesar da vantagem de recuperar energia atualmente desprezada, subsiste como principal obstáculo técnico a complexidade de integração no veículo.
\end{itemize}
A análise conjunta destas áreas evidencia que as tecnologias emergentes já oferecem respostas promissoras para as fragilidades do automóvel elétrico atual, mas nenhuma resolve isoladamente todos os requisitos. Persistem barreiras específicas --- custos iniciais de produção, densidade energética limitada, consumo dos atuadores, calibração de IA e complexidade de integração --- que justificam a abordagem sistémica deste projeto: conjugar armazenamento avançado, controlo inteligente e amortecimento regenerativo para viabilizar os veículos das próximas décadas.
\end{campotexto}

\subsection{Objetivos}
\begin{campotexto}
\instrucao{Utilizar objetivos SMART (máx. 1000 caracteres):}
\begin{itemize}[leftmargin=1.4em, itemsep=2pt, topsep=4pt]
  \item \instrucao{\textbf{Specific (Específicos):} O objetivo deve ser claro e bem definido. O que exatamente se pretende alcançar?}
  \item \instrucao{\textbf{Measurable (Mensuráveis):} Deve ser possível medir o seu progresso e saber quando o objetivo foi atingido.}
  \item \instrucao{\textbf{Achievable (Alcançáveis):} O objetivo deve ser realista e alcançável, considerando os recursos e restrições disponíveis.}
  \item \instrucao{\textbf{Relevant (Relevantes):} O objetivo deve ser importante para o projeto e estar alinhado com as metas de longo prazo.}
  \item \instrucao{\textbf{Time-bound (Temporais):} O objetivo deve ter um prazo definido, com uma data de início e fim.}
\end{itemize}
\instrucao{Esta metodologia ajuda a transformar ideias vagas em planos de ação concretos e focados.}
\end{campotexto}

\subsection{Interdisciplinaridade}
\begin{campotexto}
\instrucao{Aspetos interdisciplinares para a integração de conhecimentos, métodos e conceitos de múltiplas disciplinas de engenharia e não-engenharia que permitam resolver problemas complexos e criar soluções mais inovadoras e completas. (Máx. 1000 caracteres)}
\end{campotexto}
